% Generated draft: requires mathematical and source review.
\documentclass[11pt]{article}
\usepackage[margin=25mm]{geometry}
\usepackage{fontspec}
\setmainfont{Latin Modern Roman}
\usepackage{amsmath,amssymb,mathtools,unicode-math}
\setmathfont{Latin Modern Math}
\usepackage{enumitem,needspace,xurl,hyperref}
\setlength{\parindent}{0pt}
\setlength{\parskip}{5pt}
\setlength{\emergencystretch}{3em}
\newcommand{\N}{\mathbb{N}}
\newcommand{\Z}{\mathbb{Z}}
\newcommand{\Q}{\mathbb{Q}}
\newcommand{\R}{\mathbb{R}}
\newcommand{\C}{\mathbb{C}}
\newcommand{\F}{\mathbb{F}}
\newcommand{\A}{\mathbb{A}}
\newcommand{\PP}{\mathbb{P}}
\newcommand{\E}{\mathbb{E}}
\newcommand{\eps}{\varepsilon}
\newcommand{\dd}{\,\mathrm{d}}
\begin{document}
\section*{UnsolvedMath ID: 20000452}
\textbf{Source problem code:} catalogue entry
\subsection{Definitions and mathematical statement}
Ellenberg's AIM question assumes that for every Zariski-dense open \(U\subset X\), the arithmetic \'etale fundamental-group map admits a section. In the geometrically connected setting, with suitable geometric base points, this means
\[
\forall U\ \exists s_U:G_k\to\pi_1(U),\qquad
\pi_1(U)\to G_k\to\pi_1(U)\text{ is split by }s_U,
\]
more explicitly the structural map composed with \(s_U\) is \(\operatorname{id}_{G_k}\), and the section is continuous. Different \(U\) may have unrelated sections; a compatible family is not assumed. The question is whether
\[
X(\mathbb A_k)\ne\varnothing\quad\Longrightarrow\quad
X(\mathbb A_k)^{\operatorname{Br}}\ne\varnothing
\]
under this hypothesis. The original says only a variety over a number field, leaving smoothness, geometric connectedness, and the precise fundamental-group/adelic setting unstated. It notes Zariski-dense \(X(k)\) as a sufficient condition for the splitting hypothesis.

Poonen's remark asks whether the splitting assumption itself forces adelic solubility; McCallum proposes the corresponding local-field question, and Ellenberg asks for analogous conclusions about other obstructions. Those refinements are not answered by simply assuming the displayed adelic premise.

The attached report proposes a curve reduction. For every smooth projective geometrically connected curve \(C/k\), it asserts that finitely many closed points can be removed to obtain an affine open \(U_0\) with \(\operatorname{Pic}(U_0)=0\). A section of \(\pi_1(U_0)\to G_k\) then forces
\[
\operatorname{Br}(k)\hookrightarrow\operatorname{Br}(C).
\]
Thus the all-open splitting hypothesis is said to kill the constant relative Brauer kernel, \(\ker(\operatorname{Br}(k)\to\operatorname{Br}(C))=0\). In genus zero the report concludes \(C(k)\ne\varnothing\), without a prior adelic assumption. Injectivity of constants is not itself asserted to compute the whole Brauer--Manin set for every positive-genus curve.

For proper hyperbolic curves, the report gives a further sufficient compact-compatible selection condition on open-by-open sections that promotes them to a birational section of the function-field Galois projection. It then claims a Brauer-orthogonal adelic point. The summary does not specify the topology, compact subsets, transition maps, or exact finite-compatibility hypotheses needed for that inverse-limit argument, nor all hypotheses of the birational-section implication. These are additional conditions, not consequences supplied by \(\forall U\exists s_U\) alone.

A separate Kummer calculation on \(\mathbb G_m\) is said to exhibit ramification escape and show that required compactness cannot formally be deduced from nonempty section sets. The actual calculation is absent. Hence the report isolates a compatibility obstacle for positive genus rather than solving the original all-open question for all varieties.
\subsection{Short English statement}
Does having an arithmetic fundamental-group section on every dense open exclude a Brauer--Manin obstruction? The report gives a curve Brauer-kernel reduction and a genus-zero consequence, while recording the extra compactness and compatibility needed to pass to a birational section.
\subsection{Sources}
\begin{itemize}
\item AIM workshop ``Rational and integral points on higher dimensional varieties,'' Question 19, attributed to Ellenberg, with remarks by Poonen and McCallum: \url{https://aimath.org/WWN/qptsurface2/qptsurface2.pdf}. No PDF was accessed here.
\item Local canonical registry \texttt{aim-algebraic-number-theory-notes.json}, \texttt{notes[103]}, tag \texttt{question}; permanent UnsolvedMath ID \(20000452\), source code \texttt{catalogue entry}. The original contains damaged \(\overline k\) and \(\ne\) typography, and its nested remarks ask about adelic solubility, local fields, and other obstructions. Original literature field is empty.
\item Supplied research attempt 1 proposing the Picard-trivial affine-open/Brauer-kernel theorem, genus-zero consequence, and compact-compatible selection audit. No authoritative full references or explicit Kummer calculation are supplied. The report is machine generated, not peer reviewed, and requires independent verification before citation.
\end{itemize}
\end{document}
